\documentclass[10pt,a4paper]{amsart}
\usepackage{amsmath,amssymb,amsthm,graphicx}
% For commutative diagrams we use Xy-pic
\usepackage[all]{xy}

% We use 2cell for 2-commutative diagrams.
\xyoption{2cell}
% Theorem environments.
%
\theoremstyle{plain}
\newtheorem{theorem}[subsection]{Theorem}
\newtheorem{proposition}[subsection]{Proposition}
\newtheorem{lemma}[subsection]{Lemma}

\theoremstyle{definition}
\newtheorem{definition}[subsection]{Definition}
\newtheorem{example}[subsection]{Example}
\newtheorem{exercise}[subsection]{Exercise}
\newtheorem{situation}[subsection]{Situation}

\theoremstyle{remark}
\newtheorem{remark}[subsection]{Remark}
\newtheorem{remarks}[subsection]{Remarks}

\numberwithin{equation}{subsection}

% Macros
%
\def\lim{\mathop{\mathrm{lim}}\nolimits}
\def\colim{\mathop{\mathrm{colim}}\nolimits}
\def\Spec{\mathop{\mathrm{Spec}}}
\def\Hom{\mathop{\mathrm{Hom}}\nolimits}
\def\Ext{\mathop{\mathrm{Ext}}\nolimits}
\def\SheafHom{\mathop{\mathcal{H}\!\mathit{om}}\nolimits}
\def\SheafExt{\mathop{\mathcal{E}\!\mathit{xt}}\nolimits}
\def\Sch{\mathit{Sch}}
\def\Mor{\mathop{\mathrm{Mor}}\nolimits}
\def\Ob{\mathop{\mathrm{Ob}}\nolimits}
\def\Sh{\mathop{\mathit{Sh}}\nolimits}
\def\NL{\mathop{N\!L}\nolimits}
\def\CH{\mathop{\mathrm{CH}}\nolimits}
\def\proetale{{pro\text{-}\acute{e}tale}}
\def\etale{{\acute{e}tale}}
\def\QCoh{\mathit{QCoh}}
\def\Ker{\mathop{\mathrm{Ker}}}
\def\Im{\mathop{\mathrm{Im}}}
\def\Coker{\mathop{\mathrm{Coker}}}
\def\Coim{\mathop{\mathrm{Coim}}}

% Boxtimes
%
\DeclareMathSymbol{\boxtimes}{\mathbin}{AMSa}{"02}

%
% Macros for moduli stacks/spaces
%
\def\QCohstack{\mathcal{QC}\!\mathit{oh}}
\def\Cohstack{\mathcal{C}\!\mathit{oh}}
\def\Spacesstack{\mathcal{S}\!\mathit{paces}}
\def\Quotfunctor{\mathrm{Quot}}
\def\Hilbfunctor{\mathrm{Hilb}}
\def\Curvesstack{\mathcal{C}\!\mathit{urves}}
\def\Polarizedstack{\mathcal{P}\!\mathit{olarized}}
\def\Complexesstack{\mathcal{C}\!\mathit{omplexes}}
% \Pic is the operator that assigns to X its picard group, usage \Pic(X)
% \Picardstack_{X/B} denotes the Picard stack of X over B
% \Picardfunctor_{X/B} denotes the Picard functor of X over B
\def\Pic{\mathop{\mathrm{Pic}}\nolimits}
\def\Picardstack{\mathcal{P}\!\mathit{ic}}
\def\Picardfunctor{\mathrm{Pic}}
\def\Deformationcategory{\mathcal{D}\!\mathit{ef}}
\usepackage{cleveref}
\crefname{originallemma}{original lemma}{original lemmas}
\Crefname{originallemma}{Original lemma}{Original lemmas}
\makeatletter
\newlabel{lemma-cohomology-Gm-function-field-curve}{{03RG}{source}{}{sourcepointer-1}{}}
\newlabel{lemma-cohomology-Gm-function-field-curve@cref}{{[originallemma][0][]03RG}{[1][0][]source}{}{}{}}
\makeatother
\title{Vanishing of the higher direct images of the multiplicative group along the inclusion of the generic point}
\author{The Stacks Project Authors}
\date{stacks-project:etale-cohomology.tex:local-snapshot-20261011; GNU FDL 1.2 or later}
\begin{document}
\maketitle
\noindent\textit{Editorial note.} All delivered mathematical text is the author's own, copied from the pinned Stacks Project chapter file etale-cohomology.tex; no statement or proof step was rewritten, reordered or removed.\par\medskip
\noindent\textit{Editorial note.} The proof is the author's complete argument for both cases; the author notes that the second case is actually superfluous, and it is retained verbatim inside the counted span.\par\medskip
\noindent\textit{Editorial note.} Every labelled reference occurring inside the counted spans is carried in assembly.external\_references as a non-load-bearing navigation pointer (same-chapter ones at their definition site, cross-chapter ones at the invocation inside the counted span), so no load-bearing mathematics is replaced by a stub. No bibliography citation occurs inside the counted spans.\par\medskip
\noindent\textit{Editorial note.} The checked chapter list for same-name collisions is the set of stacks chapters staged in this repository: chow, cohomology, divisors, etale-cohomology, flat, more-algebra, more-morphisms, morphisms, spaces-more-morphisms, varieties together with this card's chapter etale-cohomology.\par\medskip
\noindent\textit{Editorial note.} The preamble is the chapter's own native preamble excerpt already used for this source: the Xy-pic setup (native preamble lines 14-18) and the native theorem declarations with the mathematical macros (native preamble lines 154-217).\par\medskip
\noindent\textit{Allowed standard prerequisites (editorial).}\begin{itemize}
\item The standing set-up: a smooth curve X over an algebraically closed field with the inclusions of closed points and of the generic point. Same chapter, the set-up paragraph delivered verbatim as a context span in this card (source lines 10937-10941). Supplies the morphism j and the sheaf G\_\{m,eta\}.
\item Stalks of the higher direct images at geometric points and their description by cohomology over strict henselizations. Same chapter, the description of the stalk of R\textasciicircum{}q j\_* at a geometric point by the cohomology of the pullback over the strict henselization of the base at that point. Supplies the computation on which the two cases of the proof rest.
\item Vanishing of the cohomology of the multiplicative group over a field of transcendence degree one. Same chapter, the vanishing of H\textasciicircum{}q(Spec(K), G\_m) for K of transcendence degree one over an algebraically closed field, for q at least one. Supplies the vanishing in the case of a closed point.
\end{itemize}\medskip
\ifdefined\hypertarget\hypertarget{sourcepointer-1}{}\fi
\noindent\textit{Source reference.} Original lemma 03RG: sources/stacks-etale-cohomology/pinned-original-source.tex, source lines 10907--10907. Navigation pointer invoked at proof line 11017; the counted span contains the pointer only and the cited result is not reproduced inside it, so the pointer is registered as a non-load-bearing navigation stub in the established form for Stacks material. No bibliography citation is involved.\par\medskip
\noindent
In this section, we fix an algebraically closed field $k$ and a smooth curve
$X$ over $k$. We denote $i_x : x \hookrightarrow X$ the inclusion of a closed
point of $X$ and $j : \eta \hookrightarrow X$ the inclusion of the generic
point. We also denote $X_0$ the set of closed points of $X$.

\begin{lemma}
\label{lemma-higher-direct-jstar-Gm}
For any $q \geq 1$, $R^q j_*\mathbf{G}_{m, \eta} = 0$.
\end{lemma}

\begin{proof}
We need to show that $(R^q j_*\mathbf{G}_{m, \eta})_{\bar x} = 0$ for every
geometric point $\bar x$ of $X$.

\medskip\noindent
Assume that $\bar x$ lies over a closed point $x$ of $X$. Let $\Spec(A)$
be an affine open neighbourhood of $x$ in $X$, and $K$ the fraction field
of $A$. Then
$$
\Spec(\mathcal{O}^{sh}_{X, \bar x}) \times_X \eta =
\Spec(\mathcal{O}^{sh}_{X, \bar x} \otimes_A K).
$$
The ring $\mathcal{O}^{sh}_{X, \bar x} \otimes_A K$ is a localization of
the discrete valuation ring $\mathcal{O}^{sh}_{X, \bar x}$, so it is either
$\mathcal{O}^{sh}_{X, \bar x}$ again, or its fraction field
$K^{sh}_{\bar x}$. But since some local uniformizer gets inverted, it must
be the latter. Hence
$$
(R^q j_*\mathbf{G}_{m, \eta})_{(X, \bar x)} = H_\etale^q(\Spec
K^{sh}_{\bar x}, \mathbf{G}_m).
$$
Now recall that $\mathcal{O}^{sh}_{X, \bar x} =
\colim_{(U, \bar u) \to \bar x} \mathcal{O}(U) = \colim_{A \subset B} B$
where $A \to B$ is \'etale, hence $K^{sh}_{\bar x}$ is an algebraic
extension of $K = k(X)$, and we may apply
Lemma \ref{lemma-cohomology-Gm-function-field-curve}
to get the vanishing.

\medskip\noindent
Assume that $\bar x = \bar \eta$ lies over the generic point $\eta$ of $X$ (in
fact, this case is superfluous). Then
$\mathcal{O}^{sh}_{X, \bar \eta} = \kappa(\eta)^{sep}$ and thus
\begin{eqnarray*}
(R^q j_*\mathbf{G}_{m, \eta})_{\bar \eta}
& = &
H_\etale^q(\Spec(\kappa(\eta)^{sep}) \times_X \eta, \mathbf{G}_m) \\
& = & H_\etale^q (\Spec(\kappa(\eta)^{sep}), \mathbf{G}_m) \\
& = & 0 \ \ \text{ for } q \geq 1
\end{eqnarray*}
since the corresponding Galois group is trivial.
\end{proof}

\end{document}
